% ============================================================
% Model: 平面转子
% ============================================================

\section{平面转子}
\label{sec:planar-rotor}

\begin{tipbox}[关键词标签]
平面转子 · 周期性边界条件 · 角动量 · 转动能级 · Stark效应 · 带电棒 · 简并
\end{tipbox}

% --------------------------------------------------------
% 1. 模型定义框
% --------------------------------------------------------
\begin{modelbox}[模型：一维角坐标上的自由转动]
\textbf{物理情境}：一个刚体被约束在平面内绕固定轴转动，如分子中的基团转动、带电棒在电场中的取向运动。转动惯量为 $I$，转角为 $\varphi\in[0,2\pi)$。

\textbf{哈密顿量}：
\[
\hat H=-\frac{\hbar^2}{2I}\frac{\partial^2}{\partial\varphi^2}=\frac{\hat L_z^2}{2I},
\quad\hat L_z=-i\hbar\frac{\partial}{\partial\varphi}.
\]

\textbf{边界条件}：波函数单值性要求周期性边界条件
\[
\psi(\varphi+2\pi)=\psi(\varphi).
\]

\textbf{本征解}：
\begin{itemize}[nosep]
    \item 本征函数：$\displaystyle\psi_m(\varphi)=\frac{1}{\sqrt{2\pi}}e^{im\varphi},\quad m=0,\pm1,\pm2,\ldots$
    \item 本征能量：$\displaystyle E_m=\frac{\hbar^2m^2}{2I}$
    \item 角动量本征值：$L_z=m\hbar$
\end{itemize}

\textbf{简并情况}：
\begin{itemize}[nosep]
    \item $m=0$：非简并（$E_0=0$）
    \item $m\neq0$：二重简并（$\pm m$ 对应相同能量）
\end{itemize}
\end{modelbox}

% --------------------------------------------------------
% 2. 关键公式框
% --------------------------------------------------------
\begin{formulabox}[核心公式]
\begin{enumerate}[label=\textbf{(\arabic*)}]
    \item \textbf{含时演化}：
    \[
    \psi(\varphi,t)=\sum_{m=-\infty}^{\infty}c_m\,\frac{e^{im\varphi}}{\sqrt{2\pi}}\,e^{-iE_mt/\hbar},
    \quad c_m=\int_0^{2\pi}\frac{e^{-im\varphi}}{\sqrt{2\pi}}\psi(\varphi,0)\,d\varphi.
    \]

    \item \textbf{矩形波包展开}（$\psi(\varphi,0)=A\sin^2\varphi$）：
    \[
    \sin^2\varphi=\frac12-\frac14e^{i2\varphi}-\frac14e^{-i2\varphi}
    \quad\Rightarrow\quad
    \psi(\varphi,t)=\frac{1}{\sqrt{3\pi}}\Bigl(1-e^{-i2\hbar t/I}\cos2\varphi\Bigr).
    \]
    \texteq{仅含 $m=0$ 和 $m=\pm2$ 三个成分}

    \item \textbf{电场中的平面转子}（带电均匀棒，电偶极矩 $d$）：
    \[
    H'=-d\mathcal{E}\cos\varphi=-\frac{d\mathcal{E}}{2}\bigl(e^{i\varphi}+e^{-i\varphi}\bigr).
    \]
    \texteq{微扰连接 $|m\rangle$ 与 $|m\pm1\rangle$，选择定则 $\Delta m=\pm1$}

    \item \textbf{一级能量修正}：
    \[
    E_m^{(1)}=\langle m|H'|m\rangle=0\quad(\text{因 }\langle m|\cos\varphi|m\rangle=0).
    \]
    \texteq{一级 Stark 效应为零}

    \item \textbf{二级能量修正}：
    \[
    E_m^{(2)}=\sum_{m'\neq m}\frac{|\langle m'|H'|m\rangle|^2}{E_m-E_{m'}}
    =\frac{Id^2\mathcal{E}^2}{\hbar^2}\frac{1}{4m^2-1}.
    \]
    \texteq{$m=0$ 时 $E_0^{(2)}=-\frac{Id^2\mathcal{E}^2}{\hbar^2}$；$m\neq0$ 时修正可正可负}
\end{enumerate}
\end{formulabox}

% --------------------------------------------------------
% 3. 训练点框
% --------------------------------------------------------
\begin{drillbox}[训练点与考法变体]
\trainingpoint{1} \textbf{周期性边界条件的应用}：给定角坐标 $\varphi\in[0,2\pi)$，写出薛定谔方程、通解、边界条件约束，得到量子化 $m\in\mathbb{Z}$。

\trainingpoint{2} \textbf{简并度分析}：平面转子 $m\neq0$ 为二重简并。加电场后微扰 $H'\propto\cos\varphi$ 在 $\{|+m\rangle,|-m\rangle\}$ 子空间内矩阵元为零，故仍可用非简并微扰公式。

\trainingpoint{3} \textbf{三角函数展开}：给定初态 $\psi(\varphi,0)$ 含 $\sin^2\varphi$、$\cos^2\varphi$ 或 $\sin\varphi\cos\varphi$，用倍角公式化为 $e^{im\varphi}$ 的线性组合。

\trainingpoint{4} \textbf{Stark 效应级数判定}：平面转子一级修正为零（$\langle\cos\varphi\rangle=0$），二级修正 $\propto\mathcal{E}^2$——这是``二次 Stark 效应''。对比氢原子 $n\geq2$ 的线性 Stark 效应（因 $l$ 简并）。

\trainingpoint{5} \textbf{角动量算符的对易关系}：验证 $[\hat L_z,\hat H]=0$，故 $L_z$ 为守恒量。测量 $L_z$ 得 $m\hbar$，概率为 $|c_m|^2$。
\end{drillbox}

% --------------------------------------------------------
% 4. 解题技巧框
% --------------------------------------------------------
\begin{tipbox}[快速识别与解题技巧]
\begin{itemize}
    \item \examnote{识别标志}：题干出现``平面转子''''转动惯量''''转角''''周期性边界''''带电棒''''电场中转动''，立即定位本模型。
    \item \examnote{与无限深势阱的对比}：平面转子也是``边界条件导致量子化''，但边界是周期性的（$\varphi$ 和 $\varphi+2\pi$ 为同一点），故 $m\in\mathbb{Z}$。无限深势阱是 Dirichlet 边界（$\psi=0$），故 $n\in\mathbb{Z}^+$。
    \item \examnote{能量零点}：$E_0=0$（$m=0$ 时），基态波函数为常数 $\psi_0=1/\sqrt{2\pi}$。这与无限深势阱 $E_1>0$ 不同。
    \item \examnote{三角函数速查}：
    \begin{center}
    $\sin^2\varphi=\frac{1-\cos2\varphi}{2}=\frac12-\frac14e^{i2\varphi}-\frac14e^{-i2\varphi}$\\[2pt]
    $\cos^2\varphi=\frac{1+\cos2\varphi}{2}=\frac12+\frac14e^{i2\varphi}+\frac14e^{-i2\varphi}$\\[2pt]
    $\sin\varphi\cos\varphi=\frac{\sin2\varphi}{2}=\frac{1}{4i}e^{i2\varphi}-\frac{1}{4i}e^{-i2\varphi}$
    \end{center}
    \item \examnote{Stark 效应判定}：若未微扰态有确定宇称且微扰为奇宇称，则一级修正为零。平面转子 $m\to-m$ 对应空间反演（$\varphi\to-\varphi$），$\cos\varphi$ 为偶函数，但矩阵元 $\langle m|\cos\varphi|m\rangle\propto\delta_{m,m}=0$？实际上 $\langle m|\cos\varphi|m\rangle=0$ 因为 $\int e^{-im\varphi}\cos\varphi e^{im\varphi}d\varphi=\int\cos\varphi d\varphi=0$。
\end{itemize}
\end{tipbox}

% --------------------------------------------------------
% 5. 易错点框
% --------------------------------------------------------
\begin{errorbox}{常见失误与陷阱}
\begin{itemize}
    \item \textbf{$m$ 的范围}：$m=0,\pm1,\pm2,\ldots$，不是 $m=1,2,3,\ldots$。负 $m$ 对应反向转动，与正 $m$ 能量简并。
    \item \textbf{归一化积分}：$\int_0^{2\pi}|e^{im\varphi}|^2d\varphi=2\pi$，故归一化常数为 $1/\sqrt{2\pi}$。不要误用 $1/\sqrt{\pi}$ 或 $1/\sqrt{a}$（势阱的归一化）。
    \item \textbf{含时演化中的交叉项}：$|\psi(\varphi,t)|^2$ 含 $\cos[(E_m-E_{m'})t/\hbar]$ 的振荡项。对于平面转子，$E_m\propto m^2$，故 $E_2-E_0=4E_1$ 而 $E_1-E_0=E_1$，振荡频率不是等间距的。
    \item \textbf{电场的方向}：若电场沿 $x$ 轴，$H'=-d\mathcal{E}\cos\varphi$；若沿 $y$ 轴，$H'=-d\mathcal{E}\sin\varphi$。不同方向的选择定则不同（$\sin\varphi$ 连接 $m$ 和 $m\pm1$ 但有不同相位）。
    \item \textbf{简并微扰的误判}：$m\neq0$ 为二重简并，但 $H'\propto\cos\varphi$ 在 $\{|m\rangle,|-m\rangle\}$ 子空间内的矩阵元为零，故可用非简并公式。若微扰为 $\sin\varphi$，则矩阵元可能非零，需严格对角化。
\end{itemize}
\end{errorbox}

% --------------------------------------------------------
% 6. 物理图像与关联模型
% --------------------------------------------------------
\begin{insightbox}[物理图像与模型关联]
\textbf{物理图像}：

平面转子就像一个在圆形轨道上跑步的人：
\begin{itemize}[nosep]
    \item 如果他不跑（$m=0$），站在原地，能量为零。
    \item 如果他顺时针跑（$m>0$）或逆时针跑（$m<0$），跑得越快（$|m|$ 越大），动能越大（$E\propto m^2$）。
    \item 顺时针和逆时针跑得一样快时，能量相同——这就是二重简并。
    \item 如果在轨道上施加一个``偏向力''（电场），跑步的人会稍微偏离原来的速度，但平均位置不变（一级修正为零），只是能量有微小变化（二级修正）。
\end{itemize}

\textbf{与相似模型的对比}：
\begin{center}
\begin{tabular}{llll}
\toprule
\textbf{对比维度} & \textbf{平面转子} & \textbf{3D 刚性转子} & \textbf{无限深势阱} \\
\midrule
坐标 & $\varphi\in[0,2\pi)$ & $(\theta,\varphi)$ 球面 & $x\in[0,a]$ \\
波函数 & $e^{im\varphi}/\sqrt{2\pi}$ & $Y_{jm}(\theta,\varphi)$ & $\sin(n\pi x/a)$ \\
能级 & $E_m\propto m^2$ & $E_j\propto j(j+1)$ & $E_n\propto n^2$ \\
简并 & $|m|$ 二重 & $2j+1$ 重 & 无 \\
Stark效应 & 二级 & 二级 & 不适用 \\
\bottomrule
\end{tabular}
\end{center}

\textbf{推广方向}：
\begin{itemize}[nosep]
    \item 双转子耦合：两个平面转子的相互作用，与自旋链模型类比
    \item 磁场中的带电转子：Aharonov-Bohm 效应在环上的实现
    \item 分子马达：生物分子（如 ATP 合成酶）的转动可用量子转子模型描述
\end{itemize}
\end{insightbox}

% --------------------------------------------------------
% 7. 推导附录
% --------------------------------------------------------
\begin{derivationbox}[推导细节（自我检验用）]
\textbf{Step 1：本征方程求解}

$-\frac{\hbar^2}{2I}\frac{d^2\psi}{d\varphi^2}=E\psi$。

令 $m^2=2IE/\hbar^2$，解为 $\psi(\varphi)=e^{im\varphi}$。

周期性边界条件 $\psi(\varphi+2\pi)=\psi(\varphi)$：$e^{im(\varphi+2\pi)}=e^{im\varphi}$ $\Rightarrow$ $e^{i2\pi m}=1$ $\Rightarrow$ $m\in\mathbb{Z}$。

归一化：$\int_0^{2\pi}|\psi_m|^2d\varphi=2\pi|c|^2=1$ $\Rightarrow$ $c=1/\sqrt{2\pi}$。

\textbf{Step 2：二级 Stark 修正推导}

$H'=-d\mathcal{E}\cos\varphi=-\frac{d\mathcal{E}}{2}(e^{i\varphi}+e^{-i\varphi})$。

矩阵元：
\[
\langle m'|H'|m\rangle=-\frac{d\mathcal{E}}{2}\bigl(\delta_{m',m+1}+\delta_{m',m-1}\bigr).
\]

对 $m\neq0$：
\begin{align*}
E_m^{(2)}&=\sum_{m'\neq m}\frac{|\langle m'|H'|m\rangle|^2}{E_m-E_{m'}}\\
&=\frac{d^2\mathcal{E}^2}{4}\Bigl(\frac{1}{E_m-E_{m+1}}+\frac{1}{E_m-E_{m-1}}\Bigr)\\
&=\frac{d^2\mathcal{E}^2}{4}\cdot\frac{2I}{\hbar^2}\Bigl(\frac{1}{m^2-(m+1)^2}+\frac{1}{m^2-(m-1)^2}\Bigr)\\
&=\frac{Id^2\mathcal{E}^2}{2\hbar^2}\Bigl(\frac{1}{-2m-1}+\frac{1}{2m-1}\Bigr)\\
&=\frac{Id^2\mathcal{E}^2}{2\hbar^2}\cdot\frac{2}{4m^2-1}
=\frac{Id^2\mathcal{E}^2}{\hbar^2}\frac{1}{4m^2-1}.
\end{align*}

对 $m=0$：只有 $m'=\pm1$ 贡献：
\[
E_0^{(2)}=2\cdot\frac{d^2\mathcal{E}^2}{4}\cdot\frac{1}{0-E_1}=-\frac{d^2\mathcal{E}^2}{2}\cdot\frac{2I}{\hbar^2}=-\frac{Id^2\mathcal{E}^2}{\hbar^2}.
\]
\end{derivationbox}

\vspace{1em}
